\documentclass[%
	% reprint,
	preprint, % !!!
	% preprintnumbers,
	amsmath,
	amssymb,
	aps,
	prx,
	floatfix,
	11pt % !!!
]{revtex4-2}
\usepackage{siunitx}
\begin{document}

\section{Network model}


We model a spatially embedded recurrent spiking neural network consisting of excitatory (E) and inhibitory (I) neuronal populations with conductance-based synapses and spike-frequency adaptation in the excitatory population.


\subsection{Single neuron dynamics}

Both excitatory and inhibitory neurons follow a leaky integrate-and-fire model with spike-frequency adaptation in the excitatory population~\cite{Treves1993}. The membrane potential $V_i$ of neuron $i$ evolves according to:
\begin{equation}
    C \frac{dV_i}{dt} = -g_L(V_i - V_L) - g_{K,i}(V_i - V_K) + I_{i},
\end{equation}
where $C$ is the membrane capacitance, $g_L$ is the leak conductance, $V_L$ is the leak reversal potential, $g_{K,i}$ is the spike-triggered potassium conductance providing adaptation, $V_K$ is the potassium reversal potential, and $I_{i}$ is the total synaptic input current.

The adaptation conductance $g_{K,i}$ decays exponentially with time constant $\tau_K$ and is incremented by $\Delta g_K$ following each spike:
\begin{equation}
    \frac{dg_{K,i}}{dt} = -\frac{g_{K,i}}{\tau_K} + \Delta g_K \sum_k \delta(t - t_i^{(k)}),
\end{equation}
where $t_i^{(k)}$ denotes the $k$-th spike time of neuron $i$. For inhibitory neurons, $\Delta g_K = 0$, so they exhibit no spike-frequency adaptation.

When the membrane potential reaches the threshold $V_{\mathrm{th}}$, the neuron emits a spike, the membrane potential is reset to $V_{\mathrm{rt}}$, and the neuron enters an absolute refractory period of duration $\tau_{\mathrm{ref}}$ during which the membrane potential is clamped. Membrane potentials are initialized by drawing from a uniform distribution $V_i(0) \sim \mathcal{U}(V_{\mathrm{rt}}, V_{\mathrm{th}})$.

\subsection{Synaptic dynamics}

Synaptic transmission is modeled using conductance-based synapses with dual-exponential kinetics. For a synapse from presynaptic neuron $j$ to postsynaptic neuron $i$, the synaptic conductance $g_{ij}(t)$ evolves according to a pair of coupled differential equations:
\begin{align}
    \frac{dg_{ij}}{dt} &= -\frac{g_{ij}}{\tau_d} + h_{ij}, \\
    \frac{dh_{ij}}{dt} &= -\frac{h_{ij}}{\tau_r} + A \left(\frac{1}{\tau_r} - \frac{1}{\tau_d}\right) \bar{g}_{ij} \sum_k \delta(t - t_j^{(k)} - d_{ij}),
\end{align}
where $h_{ij}$ is an auxiliary variable governing the rise phase, $\bar{g}_{ij}$ is the maximal synaptic conductance, $t_j^{(k)}$ is the $k$-th spike time of the presynaptic neuron, $d_{ij}$ is the synaptic transmission delay, $\tau_r$ is the rise time constant, and $\tau_d$ is the decay time constant. The normalization factor
\begin{equation}
    A = \frac{\tau_d}{\tau_d - \tau_r} \left(\frac{\tau_r}{\tau_d}\right)^{\tau_r/(\tau_r - \tau_d)}
\end{equation}
ensures that the peak conductance following a single spike equals $\bar{g}_{ij}$. The rise times are drawn from a normal distribution $\tau_r \sim \mathcal{N}(\bar{\tau}_r, 0.05\bar{\tau}_r)$ to introduce variability in effective transmission delays.

The total synaptic current to neuron $i$ is given by:
\begin{equation}
    I_{i} = -\sum_{j \in E} g_{ij}^E(t)(V_i - V_E) - \sum_{j \in I} g_{ij}^I(t)(V_i - V_I),
\end{equation}
where the sums run over all excitatory (E) and inhibitory (I) presynaptic partners, and $V_E$ and $V_I$ are the reversal potentials for excitatory and inhibitory synapses, respectively.

\subsection{Connectivity}

Neurons are distributed on a two-dimensional square domain of side length $L$ with periodic boundary conditions.
Excitatory neurons are positioned on a regular grid with density $\rho$, yielding $N_E = \rho A$ excitatory neurons arranged on an $n_E \times n_E$ grid where $n_E = \sqrt{\rho A}$.
Inhibitory neurons are uniformly distributed at random positions with density $\rho/\gamma$, giving $N_I = N_E/\gamma$ inhibitory neurons, where $\gamma$ is the population excitatory-to-inhibitory (E:I) ratio.

Synaptic connections are established probabilistically based on the Euclidean distance between neurons, with periodic boundary conditions. The connection probability from neuron $j$ at position $\mathbf{x}_j$ to neuron $i$ at position $\mathbf{x}_i$ follows an exponential kernel:
\begin{equation}
    p_{ij} = p_{\max}^{XY} \exp\left(-\frac{r_{ij}}{\sigma_{XY}}\right),
\end{equation}
where $r_{ij} = \|\mathbf{x}_i - \mathbf{x}_j\|_{\mathrm{torus}}$ is the toroidal distance between neurons, $X,Y \in \{E, I\}$ denote the pre- and postsynaptic population types, $\sigma_{XY}$ is the spatial scale of connectivity, and $p_{\max}^{XY}$ is the maximum connection probability.

Rather than specifying $p_{\max}^{XY}$ directly, we parameterize the connectivity by the mean number of incoming connections $K_{XY}$ per postsynaptic neuron. The total mass of the probability kernel over the domain is:
\begin{equation}
    \omega_{XY} = 2\pi \sigma_{XY}^2 p_{\max}^{XY},
\end{equation}
from which the expected in-degree is $\langle k_{XY} \rangle = \rho_X \omega_{XY}$, where $\rho_X$ is the density of the presynaptic population. Setting $\langle k_{XY} \rangle = K_{XY}$ determines:
\begin{equation}
    p_{\max}^{XY} = \frac{K_{XY}}{2\pi \sigma_{XY}^2 \rho_X}.
\end{equation}
The total number of synaptic connections in projection $X \to Y$ is then $K_{XY} N_Y$, and these are sampled using weighted Gumbel-max sampling to ensure the exact number of connections while respecting the distance-dependent probabilities~\cite{Kool2019}.

\subsection{Synaptic Weights}

Synaptic weights are scaled to ensure that recurrent dynamics remain in a balanced regime~\cite{}. The mean weight of synapses onto postsynaptic neuron $i$ is inversely proportional to the square root of its in-degree $k_i$:
\begin{equation}
    \langle w_{ij} \rangle_j = \frac{J^{XY}_{\mathrm{rec}}}{\sqrt{k_i}},
\end{equation}
where the constant $J^{XY}_{\mathrm{rec}}$ is computed to ensure that the population-averaged total synaptic weight equals the target weight $J_{XY}$:
\begin{equation}
    J^{XY}_{\mathrm{rec}} = J_{XY} \frac{\sum_i k_i}{\sum_i \sqrt{k_i}}.
\end{equation}
Individual synaptic weights are then drawn from a narrow normal distribution $w_{ij} \sim \mathcal{N}(\langle w_{ij} \rangle, 0.05\langle w_{ij} \rangle)$ to introduce heterogeneity while preserving the mean.

The baseline synaptic weights for excitatory projections are $J_{EE}$ and $J_{EI}$. Inhibitory synaptic weights are parameterized through a synaptic I:E ratio $\delta$:
\begin{equation}
    J_{IE} = \delta \cdot J_{EE}, \quad J_{II} = \delta \cdot J_{EI}.
\end{equation}
This parameterization allows systematic exploration of the excitation-inhibition balance by varying a single parameter $\delta$.

\subsection{External input}

Each neuron receives external excitatory input from an independent Poisson process. An external population of $N_{\mathrm{ext}}$ Poisson neurons fires at rate $\nu_{\mathrm{ext}}$, with each external neuron connected to neurons in both populations with probability $p_{\mathrm{ext}} = \sqrt{n_{\mathrm{ext}}/N_E}$, where $n_{\mathrm{ext}}$ is the target number of external synapses per excitatory neuron. External synapses use the same dual-exponential kinetics as recurrent excitatory synapses, with synaptic weights $J_{EE}$ for external-to-excitatory connections and $J_{EI}$ for external-to-inhibitory connections.

\subsection{Model parameters}

Table~\ref{tab:params} summarises the default parameter values used in simulations. The membrane parameters for excitatory and inhibitory neurons differ primarily in the leak conductance, with $g_L^E = 0.0167$~$\mu$S for excitatory neurons and $g_L^I = 0.025$~$\mu$S for inhibitory neurons, reflecting the higher intrinsic excitability of inhibitory interneurons~\cite{}.

\begin{table}[h]
\centering
\caption{Default model parameters.}
\label{tab:params}
\begin{tabular}{lll}
\hline
	\textbf{Parameter} & \textbf{Symbol} & \textbf{Value} \\
\hline
\multicolumn{3}{l}{\textit{Network geometry}} \\
Excitatory density & $\rho$ & \SI{20000}{\per\square\milli\metre} \\
Domain side length & $L$ & \SI{0.5}{\milli\metre} \\
Population E:I ratio & $\gamma$ & \num{4} \\
\hline
\multicolumn{3}{l}{\textit{Membrane parameters}} \\
Membrane capacitance & $C$ & \SI{0.25}{\nano\farad} \\
Leak conductance (E) & $g_L^E$ & \SI{0.0167}{\micro\siemens} \\
Leak conductance (I) & $g_L^I$ & \SI{0.025}{\micro\siemens} \\
Leak reversal & $V_L$ & \SI{-70}{\milli\volt} \\
Potassium reversal & $V_K$ & \SI{-85}{\milli\volt} \\
Spike threshold & $V_{\mathrm{th}}$ & \SI{-50}{\milli\volt} \\
Reset potential & $V_{\mathrm{rt}}$ & \SI{-70}{\milli\volt} \\
Refractory period & $\tau_{\mathrm{ref}}$ & \SI{4}{\milli\second} \\
Adaptation time constant & $\tau_K$ & \SI{40}{\milli\second} \\
Adaptation strength & $\Delta g_K$ & \SI{0.002}{\micro\siemens} \\
\hline
\multicolumn{3}{l}{\textit{Synaptic parameters}} \\
Excitatory rise time & $\tau_r^E$ & \SI{1.0}{\milli\second} \\
Inhibitory rise time & $\tau_r^I$ & \SI{2.0}{\milli\second} \\
Excitatory decay time & $\tau_d^E$ & \SI{5.0}{\milli\second} \\
Inhibitory decay time & $\tau_d^I$ & \SI{4.5}{\milli\second} \\
Excitatory reversal & $V_E$ & \SI{0}{\milli\volt} \\
Inhibitory reversal & $V_I$ & \SI{-80}{\milli\volt} \\
Excitatory delay & $d_E$ & \SI{1.5}{\milli\second} \\
Inhibitory delay & $d_I$ & \SI{1.5}{\milli\second} \\
\hline
\multicolumn{3}{l}{\textit{Connectivity}} \\
E$\to$E radius & $\sigma_{EE}$ & \SI{0.06}{\milli\metre} \\
E$\to$I radius & $\sigma_{EI}$ & \SI{0.07}{\milli\metre} \\
I$\to$E radius & $\sigma_{IE}$ & \SI{0.14}{\milli\metre} \\
I$\to$I radius & $\sigma_{II}$ & \SI{0.14}{\milli\metre} \\
E$\to$E in-degree & $K_{EE}$ & \num{260} \\
E$\to$I in-degree & $K_{EI}$ & \num{340} \\
I$\to$E in-degree & $K_{IE}$ & \num{225} \\
I$\to$I in-degree & $K_{II}$ & \num{290} \\
E$\to$E weight & $J_{EE}$ & \num{0.00105} \\
E$\to$I weight & $J_{EI}$ & \num{0.00145} \\
Synaptic I:E ratio & $\delta$ & \num{4.0} \\
\hline
\multicolumn{3}{l}{\textit{External input}} \\
External rate & $\nu_{\mathrm{ext}}$ & \SI{10}{\hertz} \\
External synapses & $n_{\mathrm{ext}}$ & \num{100} \\
\hline
\end{tabular}
\end{table}


\end{document}